\begin{frame}{Small to large}
    \metroset{block=fill}

    \footnotesize 

    \begin{block}{Problema}
        São dados dois inteiros $n$ e $q$, com $n, q \leq 10^5$. Considere que inicialmente temos $n$ conjuntos $\{1\}, \{2\}, \dots, \{n\}$. Serão feitas $q$ consultas, podendo ser de dois tipos:
        \begin{itemize}
            \item $1$ $i$ $j$ - faça a união do conjunto que contém $i$ com o conjunto que contém $j$.
            \item $2$ $i$ $j$ - responda se $i$ e $j$ estão no mesmo conjunto. 
        \end{itemize}
    \end{block}

    Vamos primeiro analisar o algoritmo força bruta. Declaramos um vector com $n$ sets, com o $i$-ésimo set contendo inicialmente apenas o inteiro $i$. Para cada inteiro $i$, armazenamos o índice do conjunto que o contém.

    Para cada operação de união, inserimos todos os elementos do set de índice $i$ no set de índice $j$ (e atualizamos o índice do set destes elementos), e zeramos o set $i$.

    Para cada operação do tipo $2$, apenas verificamos os índices do conjunto que as contém são iguais. Complexidade total: $O(nq \log n)$.

    % \begin{block}{Implementação da busca binária}
    %     \footnotesize
    %     \inputminted{cpp}{codigos/busca.binaria.cpp}    
    % \end{block}

\end{frame}
